\section{引言}
\label{sec:intro}

基于大语言模型的命令行编码代理（Claude Code、Codex CLI、Gemini CLI 等）已能在终端中读写代码、执行命令、运行测试并以对话推进任务。但当任务从“回答一个问题”扩展为“完成一个多步骤、跨文件、需反复编译测试的开发目标”时，它们暴露出一个共性瓶颈：\textbf{无法在无人值守下持续推进}。一次长程任务会被大量“中断点”打断——完成一步后停下等“继续”、执行敏感操作时弹确认框、面临多个方案时列菜单、遇到限流或余额不足时停在错误上。这些中断点使“让代理跑一整晚把功能做完”在实践中难以成立。

WhatyTerm 正是为消除这些中断点而设计的编排系统：它在编码代理与人类之间插入一个自动化监督层，实时感知终端状态、判断代理是否需要外部动作，并代替人类注入相应操作。系统的设计蓝图包含三项机制：规则优先、LLM 兜底的分层决策（\S\ref{sec:method-decision}）；Hook 事件、终端抓屏、进程指纹三通道感知（\S\ref{sec:arch-perception}）；以及 micro/meso/macro 三层嵌套自主循环（\S\ref{sec:method-loop}）。这些机制在设计层面自洽，也构成了本文第~\ref{sec:arch}--\ref{sec:app} 节所述的系统。

\textbf{然而，本文的核心不是宣称这套设计有效，而是报告它在真实运行中并未按预期工作。} 我们对一次真实部署的完整日志（跨越三天、覆盖四个会话、共 46 次自动决策）做了分析，得到一个与设计初衷直接冲突的结果（\S\ref{sec:reality}）：

\begin{itemize}
  \item 全部 46 次自动决策\textbf{ 100\% }由零 Token 的规则层做出，本应处理复杂场景的 LLM 兜底层\textbf{一次也未被触发}；
  \item 其中 87\%（40/46）的动作只是机械地发送“继续”，仅有 \textbf{1 次}注入了真实的任务指令；
  \item 与之呼应，系统的自主执行引擎在同期反复产出空输出与执行错误，从未真正交付内容。
\end{itemize}

也就是说，一个本应“智能编排自主开发”的系统，在真实运行中退化成了一个\textbf{“自动点继续”的宏}。这一现象并非偶然的实现 bug，而是暴露了自主代理 harness 中一个更普遍、更隐蔽的设计陷阱：\emph{当廉价的规则层被赋予对高频状态的最终决定权，而缺乏“该动作是否真正推进了任务”的反馈判断时，系统会稳定地收敛到“看起来在工作、实则空转”的局部最优}。规则层的高命中率——一个本被当作“成本优势”的指标——恰恰成了智能缺位的伪装。

基于此，本文的贡献重新定位如下：

\begin{itemize}
  \item \textbf{一个诚实的负面结果（\S\ref{sec:reality}）}。基于真实部署日志，量化并展示一个自主开发 harness 如何退化为“点继续”宏，包括决策类型分布、规则/LLM 占比与动作构成。
  \item \textbf{退化的代码级根因分析（\S\ref{sec:rootcause}）}。定位到规则层“检测到空闲提示符即无条件发送继续”的判定路径，说明为何 LLM 兜底层永不触发，并指出这是一个可推广的“规则层—推理层边界”反模式。
  \item \textbf{系统设计的完整披露（\S\ref{sec:arch}--\ref{sec:app}）}。如实呈现原始设计蓝图（分层决策、三通道感知、三层自主循环、领域策略框架），使读者能对照“设计意图”与“实际行为”的落差。
  \item \textbf{对自主 harness 设计的反思与改进方向（\S\ref{sec:discussion}）}，包括如何为规则层引入“任务推进度”反馈、如何设定规则层向推理层的强制移交条件。
\end{itemize}

本文所有定量结论均来自可查证的真实运行日志（SQLite 历史记录与引擎输出文件），文中不使用任何未经测量的臆造数据。
